\documentclass[../partie1.tex]{subfiles}
\begin{document}
%%%%%%%%%%%%%%%%%%%%%%%%%%%%%%%%%%%%%%%
% EXERCICE
\begin{exercise}[points=11]
    Résous les équations suivantes.
    \begin{enumerate}[cols=2]
        \item %4 pts
              $x^{3}+4x^{2}-3x -18=0$
        \item %2pts
              $9x^{2}+1=6x$
        \item %3 pts
              $3x^{3}-27x=0$
        \item %2 pts
              $x^{2}+2x-8 =0$
    \end{enumerate}
\end{exercise}
% SOLUTION
\begin{solution}
    \begin{enumerate}
        \item Évaluation : 
        \begin{itemize}
            \item Trouver une racine de $P(x)$\hfill 1
            \item Factoriser avec Horner\hfill 1
            \item Factoriser avec les produits remarquables\hfill 1
            \item Solutions \hfill 1 
        \end{itemize} 
        $x^{3}+4x^{2}-3x -18 = (x-2)(x+3)^2=0 \Leftrightarrow x = 2\text{~ou~} x= -3$
        
        \item Évaluation :
        \begin{itemize}
            \item Factoriser avec les produits remarquables\hfill 1
            \item Solution\hfill 1
        \end{itemize}
        $9x^2-6x+1 = (3x-1)^2 =0\Leftrightarrow x = \dfrac{1}{3}$

        \item 
        \begin{itemize}
            \item Factoriser avec la mise en évidence \hfill 1
            \item Factoriser avec les produits remarquables\hfill 1
            \item Solutions \hfill 1
        \end{itemize}
        $3x^3-27x = 3x(x-3)(x+3)=0 \Leftrightarrow x=0 \text{~ou~}x=3\text{~ou~}x=-3$
        
        \item 
        \begin{itemize}
            \item Factoriser avec somme produit\hfill 1
            \item Solution \hfill 1
        \end{itemize}
        $x^2+2x-8 = (x+4)(x-2) = 0\Leftrightarrow x=-4 \text{~ou~}x=2$
    \end{enumerate}
\end{solution}
%%%%%%%%%%%%%%%%%%%%%%%%%%%%%%%%%%%%%%%
% EXERCICE
\begin{exercise}[points=5]
    Effectue les opérations ci-dessous. Simplifie au maximum ta réponse finale et n'oublie pas les conditions d'existence.
    \begin{enumerate}[cols=2]
        \item %3 pts
              $\dfrac{4}{3x+x^2} \div \dfrac{x-3}{x^2-9}$
        \item %2 pts
              $\dfrac{x+1}{x-2}\cdot\dfrac{x^2-4x+4}{x}$
    \end{enumerate}
\end{exercise}
% SOLUTION
\begin{solution}
    \begin{enumerate}
        \item
        Évaluation :
        \begin{itemize}
            \item Factorisation \hfill 1
            \item Condition d'existence \hfill 1
            \item Simplification \hfill 0,5
        \end{itemize} 
              $\begin{aligned}[t]
                      \dfrac{4}{x(3+x)}\div\dfrac{x-3}{(x-3)(x+3)} & = \dfrac{4}{3+x}\cdot (x+3)
                      \qquad\text{ C.E. : }x\ne 3\text{ et } x\ne -3\\
                                                                   & = 4
                  \end{aligned}$
        \item
        Évaluation :
        \begin{itemize}
            \item Factorisation \hfill 1
            \item Condition d'existence \hfill 1
            \item Simplification \hfill 0,5
        \end{itemize} 
        $\begin{aligned}[t]
            \dfrac{x+1}{x-2} \cdot \dfrac{(x-2)^2}{x} 
            &= \dfrac{(x+1)(x-2)}{x}\qquad\text{C.E. : }x\ne 2 \text{ et } x \ne 0\\
        \end{aligned}$
    \end{enumerate}
\end{solution}
% EXERCICE
\begin{exercise}[points=6]
    Détermine l'expression analytique des fonctions suivantes à l'aide des informations fournies.
    \begin{enumerate}
        \item La fonction $f_1$ passe par les points $A(1;11)$ et $B(-2;2)$.
              %\item La fonction $f_2$ est linéaire et passe par le point $C(2;-8)$
        \item Le graphe de la fonction $f_2$ est le suivant. \begin{tikzpicture}[baseline=(current bounding box.east), scale=0.6]
    \tkzInit[xmin=-2,xmax=5,ymin=-3, ymax=3]
    \tkzGrid
    \tkzDrawX[right=1pt]
    \tkzLabelX[orig=false, font=\scriptsize]
    \tkzDrawY[above=1pt]
    \tkzLabelY[orig=false, font=\scriptsize]
    \tkzFct[domain=-5:5] {(3*\x -5)}
\end{tikzpicture}
        \item La racine de $f_3$ est $-2$ et son graphe est parallèle à celui de la fonction $g(x) = 2x + 2$.
              %\item Un tableau de valeur de $f_5$ est le suivant.
              %   $\begin{tblr}{hlines,vlines,columns={1cm,c}}
              %           x    & 4  & 2    & -8 \\
              %           f(x) & 13 & 12,5 & 4
              %       \end{tblr}$
    \end{enumerate}
\end{exercise}
% SOLUTION
\begin{solution}
    Évaluation :
    \begin{itemize}
        \item Pente \hfill 1
        \item Ordonnée à l'origine \hfill 1
    \end{itemize}
    \begin{enumerate}[cols=3]
        \item $f_1(x) = 3x + 8$
        \item $f_2(x) = 3x - 5$
        \item $f_3(x) = 2x + 4$
    \end{enumerate}
\end{solution}
%%%%%%%%%%%%%%%%%%%%%%%%%%%%%%%%%%%%%%%
% EXERCICE
\begin{exercise}[points=1]
    Écris les expressions ci-dessous sans dénominateur en utilisant les propriétés des puissances.
    \begin{enumerate}[cols=2]
        \item $\dfrac{15x^{-2}}{5y^7}$
        \item $\dfrac{x^7 y^{-2}}{x^{-2}y^5}$
    \end{enumerate}
\end{exercise}
% SOLUTION
\begin{solution}
    Évaluation :
    \begin{itemize}
        \item Expression correcte \hfill 0,5
    \end{itemize}
\begin{enumerate}[cols=2]
    \item $3x^{-2}y^{-7}$
    \item $x^9y^{-7}$
\end{enumerate}
\end{solution}
%%%%%%%%%%%%%%%%%%%%%%%%%%%%%%%%%%%%%%%
% EXERCICE
\begin{exercise}[points=3]
    Effectue les opérations demandées. Simplifie au maximum ta réponse.
    \begin{enumerate}[cols=3]
        \item $\sqrt{35}\cdot\sqrt{20}$
        \item $\sqrt{45}-\sqrt{80}$
        \item $\left(3 + 2\sqrt{2}\right)^2$
    \end{enumerate}
\end{exercise}
% SOLUTION
\begin{solution}
    Évaluation :
    \begin{itemize}
        \item L'expression est correcte \hfill 1
    \end{itemize}
\begin{enumerate}[cols=3]
    \item $10 \sqrt{7}$
    \item $-\sqrt{5}$
    \item $17+12\sqrt{2}$
\end{enumerate}
\end{solution}
\end{document}